%================== HIGHLIGHTS ==================
\begin{summarybox}{Highlights}
\begin{itemize}[nosep]
\item One heat-flux template, with three lags and one memory term, carries the classical, Lord--Shulman, dual-phase-lag, three-phase-lag, Moore--Gibson--Thompson, fractional and memory-dependent theories (and the Green--Naghdi types as limits) as parameter choices. Green--Lindsay and Guyer--Krumhansl lie outside it, and we treat them separately.
\item A single coupled dispersion polynomial, verified symbolically, gives comparative phase speed, attenuation and thermoelastic damping across the template models.
\item The widely repeated dual-phase-lag stability condition $\tau_T\ge\tau_q/2$ belongs to one specific Taylor truncation, not to the first-order template.
\item We derive finite thermal fronts, Rayleigh-wave damping, piezoelectric renormalisation and interface energy-partition rules, and check them numerically.
\item For journal submission we supply open problems, a pitfalls checklist, a decision guide and reporting standards.
\end{itemize}
\end{summarybox}

%================== GRAPHICAL ABSTRACT ==================
\begin{figure}[t!]
\centering
\begin{tikzpicture}[x=1cm,y=1cm,>=Stealth,
  box/.style={draw,rounded corners=1pt,fill=blue!8,align=center,font=\footnotesize,inner sep=2pt},
  arr/.style={->,thick}]
  \node[box,text width=3.2cm] (heat) {Heat flux $q_i$\\relaxation \& lags};
  \node[box,text width=3.2cm,right=0.8cm of heat] (temp) {Temperature $\theta$\\gradient};
  \node[box,text width=3.2cm,below=1.2cm of heat] (disp) {Dispersion\\branches};
  \node[box,text width=3.2cm,right=0.8cm of disp] (energy) {Energy\\partition};
  \node[box,text width=3.2cm,below=1.2cm of disp] (size) {Size effects\\nonlocal/gradient};
  \node[box,text width=3.2cm,right=0.8cm of size] (multi) {Multiphysics\\couplings};
  \draw[arr] (heat) -- (temp);
  \draw[arr] (heat.south) -- (disp.north);
  \draw[arr] (disp.south) -- (size.north);
  \draw[arr] (disp) -- (energy);
  \draw[arr] (size) -- (multi);
  \draw[arr] (multi.south) -- (energy.north);
\end{tikzpicture}
\caption{Graphical abstract: the unified template connects heat-flux relaxation to dispersion branches, size effects and multiphysics couplings, with energy partition at interfaces.}\label{fig:graphical}
\end{figure}

\begin{abstract}
A solid that conducts heat, deforms and carries further fields, such as electric polarization, magnetic induction, diffusing species or microstructural degrees of freedom, is governed by one member of a family of coupled models; the literature presents them one member at a time. This review organises the family around a single heat-flux template with three lags and one memory term. We show how the classical, Lord--Shulman, Green--Lindsay, Green--Naghdi, dual-phase-lag, three-phase-lag, Moore--Gibson--Thompson, fractional and memory-dependent theories arise, or fail to arise, as parameter choices, and we record which limits are regular and which are singular; one widely repeated stability statement for the dual-phase-lag model is corrected. For the coupled problem in a linear homogeneous medium we derive a dispersion polynomial in $K=k^2$ for the template models with rational heat laws and verify it symbolically. We state where exact delays, fractional powers and memory kernels make the relation transcendental or multivalued, and we use the polynomial to compare phase speed, attenuation and thermoelastic damping across the theories. Transient problems follow (heat pulse, thermal shock), then Rayleigh waves, reflection and interface energy partition, piezoelectric, magneto-thermoelastic, diffusive and microstructural couplings, and nonlocal and gradient size effects set against lattice dispersion. We discuss analytical and numerical solution techniques with a benchmark, applications from laser ultrasonics to micro-resonators and tissue heating, and close with a comparative synthesis, a pitfalls checklist, a decision guide and a research agenda. Flexural and plate thermoelastic damping are validated against exact solutions; we add anisotropic and free-surface reflection results, a finite-difference benchmark for hyperbolic heat conduction, worked examples and reporting standards. All quantitative figures are generated by accompanying scripts, and the key identities are checked symbolically or numerically.
\medskip

\noindent\textbf{Keywords:} generalized thermoelasticity; second sound; dual-phase-lag; Moore--Gibson--Thompson; fractional heat conduction; memory-dependent derivative; nonlocal elasticity; piezoelectricity; Rayleigh waves; thermoelastic damping; interface energy partition.
\end{abstract}
